\begingroup
\setlength{\abovedisplayskip}{6pt}
\setlength{\belowdisplayskip}{6pt}
\setlength{\abovedisplayshortskip}{4pt}
\setlength{\belowdisplayshortskip}{4pt}
\setlength{\intextsep}{6pt}

\textbf{B.8.1. Input data.} The projectile mass is denoted by $m_p$ to distinguish it from the attached mass $M$. Since the projectile remains embedded after impact, the vibrating mass is $m_T=M+m_p$. I represent the beam by its equivalent tip stiffness.

{ % do not increment counter
\begin{center}\small\begin{tabular}{@{}lrl@{}}
\toprule\noalign{}
Quantity & Value & Unit \\
\midrule\noalign{}
$L$ & 0.525 & m \\
$E$ & 200 & GPa \\
$h$ & 4.00 & mm \\
$b$ & 30.0 & mm \\
$M$ & 62.0 & kg \\
$m_p$ & 0.0130 & kg \\
$v_p$ & 1020 & m/s \\
$t_1$ & 17.0 & s \\
$N_{\mathrm{osc}}$ & 13 & oscillations \\
\bottomrule\end{tabular}\end{center}
}

\textbf{B.8.2. Second moment of area.} The beam moves horizontally, but the bending axis is determined by the cross-section orientation. In the section shown, $h$ is the bending depth and $b$ is the dimension normal to that depth. Therefore,

\[
I=\frac{bh^3}{12}
 =\frac{(0.030)(0.004)^3}{12}
 =1.600\times10^{-10}\ \mathrm{m^4}
 =160.0\ \mathrm{mm^4}.
\]

\[
\boxed{I=1.600\times10^{-10}\ \mathrm{m^4}
\qquad\left(160.0\ \mathrm{mm^4}\right)}
\qquad\text{(B.8.2)}
\]

\textbf{B.8.3. Cross-sectional stiffness.} Multiplying the elastic modulus by the second moment of area gives

\[
EI=(200\times10^9)(1.600\times10^{-10})
 =32.00\ \mathrm{N\,m^2}.
\]

\[
\boxed{EI=32.00\ \mathrm{N\,m^2}}
\qquad\text{(B.8.3)}
\]

\textbf{B.8.4. Flex-beam spring constant.} The fixed-free beam tip deflection under a concentrated force gives the equivalent spring constant

\[
k=\frac{3EI}{L^3}
 =\frac{3(32.00)}{(0.525)^3}
 =663.427\ \mathrm{N/m}.
\]

\[
\boxed{k=663.427\ \mathrm{N/m}}
\qquad\text{(B.8.4)}
\]

\textbf{B.8.5. Initial velocity after impact.} The collision duration is short compared with the vibration period, so the beam spring force contributes negligible impulse during the collision. Conservation of linear momentum therefore gives

\[
m_pv_p+M(0)=(M+m_p)v_0.
\]

Solving for the common velocity immediately after impact gives

\[
v_0=\frac{m_pv_p}{M+m_p}
 =\frac{(0.0130)(1020)}{62.013}
 =0.213826\ \mathrm{m/s}.
\]

\[
\boxed{\dot{u}_0=v_0=0.213826\ \mathrm{m/s}}
\qquad\text{(B.8.5)}
\]

\textbf{B.8.6. Natural frequency and period.} After the collision, the projectile and attached mass move together. Because the beam is horizontal and damping is neglected, the free-vibration equation is

\[
(M+m_p)\ddot{u}+ku=0.
\]

The natural frequency follows from the characteristic equation.

\[
\omega_n=\sqrt{\frac{k}{M+m_p}}
 =\sqrt{\frac{663.427}{62.013}}
 =3.27081\ \mathrm{rad/s}.
\]

The period is the time required for one complete cycle.

\[
\tau=\frac{2\pi}{\omega_n}
 =\frac{2\pi}{3.27081}
 =1.92099\ \mathrm{s}.
\]

\[
\boxed{\omega_n=3.27081\ \mathrm{rad/s},
\qquad
\tau=1.92099\ \mathrm{s}}
\qquad\text{(B.8.6)}
\]

\textbf{B.8.7. Displacement at $t_1=17\ \mathrm{s}$.} The beam displacement is continuous through the impact, so $u(0)=0$. The initial velocity is $\dot{u}(0)=v_0$. Applying these conditions to the undamped free-vibration solution gives

\[
u(t)=u(0)\cos(\omega_nt)
 +\frac{\dot{u}(0)}{\omega_n}\sin(\omega_nt)
 =\frac{v_0}{\omega_n}\sin(\omega_nt).
\]

At $t_1=17\ \mathrm{s}$,

\[
u(t_1)=\frac{0.213826}{3.27081}\sin\bigl((3.27081)(17)\bigr)
 =-0.0529817\ \mathrm{m}.
\]

\[
\boxed{u(17\ \mathrm{s})=-52.9817\ \mathrm{mm}}
\qquad\text{(B.8.7)}
\]

\textbf{B.8.8. Thirteen-oscillation response.}

\textbf{(a) Largest response.} The response amplitude is

\[
\lvert u\rvert_{\max}=\frac{v_0}{\omega_n}
 =\frac{0.213826}{3.27081}
 =0.0653741\ \mathrm{m}.
\]

The first positive maximum occurs when $\omega_nt=\pi/2$, so

\[
t_{\mathrm{peak},1}=\frac{\pi}{2\omega_n}
 =0.480247\ \mathrm{s}.
\]

\[
\boxed{\lvert u\rvert_{\max}=65.374\ \mathrm{mm}
\quad\text{at}\quad
t_{\mathrm{peak},1}=0.480247\ \mathrm{s}}
\qquad\text{(B.8.8.a)}
\]

\textbf{(b) Plot markers.} Thirteen oscillations span

\[
t_{\mathrm{end}}=N_{\mathrm{osc}}\tau
 =13(1.92099)
 =24.9728\ \mathrm{s}.
\]

The plot marks the first positive maximum and the response at $t_1=17\ \mathrm{s}$. Their coordinates are $(0.480\ \mathrm{s},65.4\ \mathrm{mm})$ and $(17.00\ \mathrm{s},-53.0\ \mathrm{mm})$.

\begin{figure}[H]
\centering
\includegraphics[width=0.95\linewidth,height=0.38\textheight,keepaspectratio,alt={B.8.8.a-b. Undamped horizontal flex-beam response after embedded-projectile impact, with markers at the first maximum and at seventeen seconds.}]{figures/generated/solutions/B8.pdf}
\caption{B.8.8.a-b. Undamped horizontal flex-beam response over $N_{\mathrm{osc}}=13$ cycles.}
\end{figure}

\textbf{B.8.9. Discussion.} The projectile transfers its momentum to the combined mass, which creates the initial velocity $v_0$ while leaving the displacement equal to zero. The undamped model then preserves the amplitude at approximately $65.4\ \mathrm{mm}$. This is about $12.5\%$ of the $525\ \mathrm{mm}$ span, so the result is large enough that small-deflection beam theory should be treated as a linear-model estimate.

The negative value at $17\ \mathrm{s}$ is caused by the phase of the sinusoidal response at that time; it does not indicate a change in amplitude. The calculation also omits distributed beam mass, damping, local impact deformation, and large-deflection geometric effects. Including those effects would change the natural frequency and reduce or otherwise modify the observed response.

\clearpage
\matlabcaption{Thirteen-oscillation response after the embedded-projectile impact.}

\lstinputlisting[style=hwmatlab]{matlab/B8.m}

\endgroup
